\section{Method}

\subsection{Data and split}
HAM10000~\cite{tschandl2018ham10000} contains \nImages{} dermatoscopic images of \nLesions{} distinct lesions in seven diagnostic classes (akiec, bcc, bkl, df, mel, nv, vasc).
Many lesions were photographed more than once, so a random image-level split would place near-duplicate images of the same lesion in both training and test sets and inflate test scores.
I therefore split by \texttt{lesion\_id} with scikit-learn's \texttt{StratifiedGroupKFold}, stratified on diagnosis, into roughly 70/15/15\% training, validation and test sets (\nTrain/\nVal/\nTest{} images; Table~\ref{tab:split}).
The split is saved to \texttt{splits/ham10000\_split.csv} and reused by every script.
The classes are highly imbalanced: melanocytic nevi (nv) make up about two thirds of the images, while dermatofibroma (df) and vascular lesions (vasc) each make up under 2\%.

\begin{table}[h]
\centering
\caption{Images per class in each split (lesion-level split).}
\label{tab:split}
\IfFileExists{results/split_table.tex}{\begin{tabular}{@{}lccc@{}}
\toprule
Class & Train & Val & Test \\ \midrule
akiec & 233 & 60 & 34 \\
bcc & 354 & 71 & 89 \\
bkl & 790 & 141 & 168 \\
df & 91 & 15 & 9 \\
mel & 779 & 188 & 146 \\
nv & 4779 & 952 & 974 \\
vasc & 100 & 21 & 21 \\
\midrule
Total & 7126 & 1448 & 1441 \\
\bottomrule
\end{tabular}
}{\textit{Generated by train\_models.py.}}
\end{table}

\subsection{Models and training}
ResNet-50 uses torchvision's ImageNet-1k weights (\texttt{IMAGENET1K\_V2}) with the final layer replaced by a 7-way linear layer.
ViT-Base/16 uses the Hugging Face \texttt{google/vit-base-patch16-224} checkpoint
with a new 7-way classification head.
Images are resized to $224\times224$ and normalised with each checkpoint's own mean and standard deviation.

Training augmentation is random resized cropping (scale 0.7--1.0), horizontal and vertical flips, rotation up to $20^\circ$ and mild colour jitter; dermatoscopic images have no canonical orientation, so these transforms preserve the label.
To counter imbalance, the cross-entropy loss is weighted per class by the inverse of its training-set frequency.

Both models are fine-tuned end to end with AdamW (weight decay 0.05), a cosine learning-rate schedule, batch size 32 and a learning rate of $10^{-4}$ for ResNet-50 and $2\times10^{-5}$ for ViT-Base, for \resnetEpochs{} epochs.
After each epoch I compute balanced accuracy on the validation set and keep the best checkpoint.
The test set is evaluated once, with that checkpoint.
No other hyperparameters were tuned.
Training ran on a MacBook Air (\hardware) with PyTorch's MPS backend.

\subsection{Metrics}
Because always predicting nv already gives about 67\% accuracy, plain accuracy says little here.
I report balanced accuracy (the mean of per-class recall), macro-averaged F1 and per-class recall, alongside accuracy.
95\% confidence intervals are percentile bootstrap intervals over test images (2,000 resamples); the difference between the two models uses a paired bootstrap on the same resampled images.
These intervals reflect test-set sampling only, not variation between training runs.
